\begin{example}[One Cohen real]
Let
\[
  \mathbb P=\operatorname{Fn}(\omega,2,{<}\omega),
\]
the finite partial functions from \(\omega\) to \(2\), ordered by reverse
inclusion.  A name is necessary here because the generic function depends on
the external filter \(G\), while membership in \(M[G]\) must be witnessed by
an object already in \(M\).  Define
\[
  \dot c\defeq
  \set{
    \tuple{\check{\tuple{n,i}},p}:
    p\in\mathbb P,\ n\in\dom(p),\ p(n)=i
  }.
\]
The symbol \(\dot c\) stands for a \(\mathbb P\)-name in \(M\), not for the
function \(c\) in the extension; the dot records precisely this type
difference.  This is nonredundant with \(\check x\): check names interpret as
fixed members of \(M\), whereas \(\dot c\) interprets differently with the generic
filter.  The ground-set-name recursion and separation in \(M\) give
\(\dot c\in M\cap V^{\mathbb P}\).  Directly from name interpretation,
\[
\begin{aligned}
  (\dot c)^G
    &=\set{\tuple{n,i}:\exists p\in G\,(p(n)=i)}\\
    &=\bigcup G.
\end{aligned}
\]
Thus, if \(c\defeq\bigcup G\), then \(c=(\dot c)^G\in M[G]\); this immediate
use is the reason for introducing \(\dot c\).  Directedness
of \(G\) makes the union single-valued, and for each \(n\) the set
\[
  D_n=\set{p\in\mathbb P:n\in\dom(p)}
\]
is dense and belongs to \(M\), so genericity makes \(c\) a total function
\(\omega\to2\).  Write \(2^\omega\) for the set of total functions
\(\omega\to2\).  For \(b\in M\cap2^\omega\), the set
\[
  E_b=\set{p\in\mathbb P:
            \exists n\in\dom(p)\,(p(n)\ne b(n))}
\]
belongs to \(M\) and is dense: extend any condition at one fresh coordinate
with the value opposite to \(b\).  Since \(G\cap E_b\ne\varnothing\),
\(c\ne b\).  Hence \(c\notin M\), while the displayed name calculation proves
that the extension contains this new real.  It does not by itself prove
\(\neg\mathsf{CH}\);
that conclusion requires a poset adding enough reals together with
cardinal-preservation arguments.
\end{example}

